\documentclass[12pt, a4paper]{article}

% --- PACKAGES ---
\usepackage[utf8]{inputenc}
\usepackage[T1]{fontenc}
\usepackage{amsmath, amssymb, amsthm}
\usepackage{mathtools}
\usepackage{physics}
\usepackage{geometry}
\usepackage{graphicx}
\usepackage{siunitx}
\usepackage{hyperref}
\usepackage{booktabs}
\usepackage{enumitem}
\usepackage{xcolor}
\usepackage{tcolorbox}
\usepackage{bm}
\usepackage{cancel}

\geometry{margin=1in}

% --- THEOREM ENVIRONMENTS ---
\theoremstyle{plain}
\newtheorem{theorem}{Theorem}[section]
\newtheorem{proposition}[theorem]{Proposition}
\newtheorem{lemma}[theorem]{Lemma}
\newtheorem{corollary}[theorem]{Corollary}

\theoremstyle{definition}
\newtheorem{definition}[theorem]{Definition}
\newtheorem{remark}[theorem]{Remark}

% --- CUSTOM COMMANDS ---
\newcommand{\vA}{\mathbf{A}}
\newcommand{\vB}{\mathbf{B}}
\newcommand{\vE}{\mathbf{E}}
\newcommand{\vJ}{\mathbf{J}}
\newcommand{\vv}{\mathbf{v}}
\newcommand{\xhat}{\hat{\mathbf{x}}}
\newcommand{\yhat}{\hat{\mathbf{y}}}
\newcommand{\zhat}{\hat{\mathbf{z}}}
\newcommand{\nhat}{\hat{\mathbf{n}}}
\newcommand{\taup}{\tau_p}
\newcommand{\oms}{\omega_s}
\newcommand{\omstar}{\omega_s^*}
\newcommand{\sig}{\sigma}
\newcommand{\muZ}{\mu_0}
\newcommand{\Rm}{R_m}
\newcommand{\Fstar}{F^*}
\newcommand{\Spar}{S}
\newcommand{\avg}[1]{\langle #1 \rangle}
\newcommand{\Der}[2]{\frac{\partial #1}{\partial #2}}
\newcommand{\DDer}[2]{\frac{\partial^2 #1}{\partial #2^2}}

% --- TITLE ---
\title{%
  \textbf{Rigorous Electrodynamic Derivation for a Planar PCB\\
  Linear Induction Motor: From Maxwell's Equations\\
  to the Closed-Form Scaling Law}\\[0.5em]
  \large Phase 2a — Complete PDE Derivation\\
  \large Solid-State EMALS Analog — Masters Thesis
}
\author{Mohana Rangan Desigan}
\date{June 2026}

% ============================================================
\begin{document}
\maketitle
\tableofcontents
\newpage

% ============================================================
\section{Introduction and Scope}

This document contains the complete, step-by-step analytical derivation for the planar
Printed-Circuit Linear Induction Motor (PCB-LIM) system. Every intermediate algebraic
step is shown explicitly. The derivation proceeds in the following sequence:

\begin{enumerate}[label=\arabic*.]
  \item Maxwell's equations in Gaussian and SI forms, specialised to the magnetoquasistatic
        (MQS) limit appropriate to the PCB-LIM frequency range.
  \item The magnetic vector potential $\vA$ in the Coulomb gauge and the derivation of
        the Convective Magnetic Diffusion Equation (CMDE) governing eddy-current
        penetration into the moving rotor.
  \item Nondimensionalisation of the CMDE, identification of the quality factor $S$ and
        magnetic Reynolds number $\Rm$ as the governing dimensionless groups, and the
        physical interpretation of each regime.
  \item The magnetic skin depth $\delta$ and the thin-sheet vs.\ skin-effect regimes with
        the transition frequency $\omega_s^{\mathrm{tr}}$.
  \item Proof of Proposition~\ref{prop:optslip} (Optimal Slip Frequency): a sinusoidal
        analysis of the boundary-value problem in the rotor half-space yields the closed-form
        scaling law $\Fstar = S/(1+S^2)$ and the optimum $S^* = 1$.
  \item Proof of Theorem~\ref{thm:volcancellation} (Volume Cancellation): Newton's law for
        the thin-sheet rotor reduces to an ODE in velocity alone, independent of rotor
        plan-area.
  \item Force calculation via the Maxwell Stress Tensor and derivation of the
        time-averaged propulsive thrust in closed form.
  \item Derivation of the kinematic ODE and its exact analytical solution.
\end{enumerate}

\begin{tcolorbox}[colback=blue!5, colframe=blue!40, title=\textbf{Compute Platform Note}]
All numerical validation of equations in this document is executed on:
\begin{itemize}
  \item \textbf{Tier 1 / 2D FEM parametric sweep:} Google Colab free tier, T4 GPU,
        15\,GB VRAM, $\approx 5$\,hours/day runtime. All FEM jobs are designed to
        complete in $\leq 90$ minutes per operating point with \texttt{.h5} checkpointing
        to Google Drive.
  \item \textbf{Tier 2 / 3D transient FEM:} Mac Mini M4, Apple M4 SoC,
        16\,GB unified RAM, macOS. Elmer FEM + Gmsh via Homebrew, OpenMP
        parallelism on 4 performance cores. No session-time limit; overnight scheduling
        used for the $\sim 8$\,hr fine-mesh runs.
\end{itemize}
\end{tcolorbox}


% ============================================================
\section{Maxwell's Equations and the Magnetoquasistatic Limit}
\label{sec:maxwell}

\subsection{Full Maxwell's Equations in SI Units}

The complete set of Maxwell's equations in a linear, isotropic medium with permeability
$\muZ$ (non-magnetic rotor, so $\mu = \muZ$) and permittivity $\varepsilon_0$ is:

\begin{align}
  \nabla \cdot \vB &= 0 \label{eq:divB}\\
  \nabla \cdot \vE &= \frac{\rho_f}{\varepsilon_0} \label{eq:divE}\\
  \nabla \times \vE &= -\Der{\vB}{t} \label{eq:faraday}\\
  \nabla \times \vB &= \muZ \vJ + \muZ \varepsilon_0 \Der{\vE}{t} \label{eq:ampere}
\end{align}

where $\rho_f$ is free charge density and $\vJ$ is the free current density.

\subsection{The Magnetoquasistatic Approximation}

The displacement current term $\muZ \varepsilon_0 \,\partial \vE / \partial t$ in
Amp\`ere's law~\eqref{eq:ampere} is negligible when the electromagnetic wavelength
$\lambda_{\mathrm{em}} = c/f$ is much larger than the system length scale $L$. For
the PCB-LIM:

\begin{align}
  f &\leq 500\,\si{Hz}, \quad L = \taup \leq 20\,\si{mm} \nonumber\\
  \lambda_{\mathrm{em}} &= \frac{3\times10^8}{500} = 600\,\si{km} \gg L \nonumber
\end{align}

The ratio $L/\lambda_{\mathrm{em}} \sim 3\times10^{-8}$, so the MQS approximation
is exact to better than one part in $10^7$. We therefore drop the displacement current
and work with:

\begin{align}
  \nabla \cdot \vB &= 0 \label{eq:mqs1}\\
  \nabla \times \vE &= -\Der{\vB}{t} \label{eq:mqs2}\\
  \nabla \times \vB &= \muZ \vJ \label{eq:mqs3}
\end{align}

\subsection{Ohm's Law in a Moving Conductor}

For a conducting body moving with velocity $\vv$ in an external magnetic field, the
Lorentz transformation of the electric field into the lab frame gives the generalized
Ohm's law:

\begin{equation}
  \vJ = \sig \left( \vE + \vv \times \vB \right)
  \label{eq:ohm}
\end{equation}

This is exact to first order in $v/c$, which is entirely adequate for $v \leq v_s \sim
10\,\si{m/s} \ll c$.


% ============================================================
\section{Magnetic Vector Potential and the Coulomb Gauge}
\label{sec:vectorpot}

\subsection{Introduction of the Vector Potential}

Since $\nabla \cdot \vB = 0$ identically, we can always write:

\begin{equation}
  \vB = \nabla \times \vA
  \label{eq:BcurlA}
\end{equation}

for some magnetic vector potential $\vA(\mathbf{r}, t)$. Substituting into
Faraday's law~\eqref{eq:mqs2}:

\begin{equation}
  \nabla \times \vE = -\frac{\partial}{\partial t}\left(\nabla \times \vA\right)
  = -\nabla \times \Der{\vA}{t}
\end{equation}

\begin{equation}
  \nabla \times \left(\vE + \Der{\vA}{t}\right) = \mathbf{0}
\end{equation}

Since the curl vanishes, the quantity in parentheses is the gradient of a scalar:

\begin{equation}
  \vE + \Der{\vA}{t} = -\nabla \phi \quad \Longrightarrow \quad
  \vE = -\nabla\phi - \Der{\vA}{t}
  \label{eq:Ephi}
\end{equation}

where $\phi$ is the electric scalar potential.

\subsection{The Coulomb Gauge}

The vector potential $\vA$ is not unique — any gradient field can be added without
changing $\vB$. We fix this gauge freedom by imposing the \emph{Coulomb gauge}:

\begin{equation}
  \nabla \cdot \vA = 0
  \label{eq:coulomb}
\end{equation}

Under this gauge, substituting $\vB = \nabla \times \vA$ into Amp\`ere's
law~\eqref{eq:mqs3} and using the vector identity
$\nabla \times (\nabla \times \vA) = \nabla(\nabla \cdot \vA) - \nabla^2 \vA$:

\begin{align}
  \muZ \vJ &= \nabla \times \vB = \nabla \times (\nabla \times \vA) \nonumber\\
            &= \nabla\underbrace{(\nabla \cdot \vA)}_{=\,0} - \nabla^2 \vA \nonumber\\
            &= -\nabla^2 \vA
\end{align}

Therefore:
\begin{equation}
  \nabla^2 \vA = -\muZ \vJ
  \label{eq:poissonA}
\end{equation}

This is the vector Poisson equation for $\vA$, valid everywhere.

% ============================================================
\section{Convective Magnetic Diffusion Equation}
\label{sec:cmde}

\subsection{Derivation Inside the Moving Conductor}

Inside the rotor (moving at velocity $\vv = v\,\xhat$), we combine the generalized
Ohm's law~\eqref{eq:ohm} with Faraday's law and the Poisson
equation~\eqref{eq:poissonA}. Substituting~\eqref{eq:Ephi} into~\eqref{eq:ohm}:

\begin{equation}
  \vJ = \sig\left(-\nabla\phi - \Der{\vA}{t} + \vv \times (\nabla \times \vA)\right)
  \label{eq:Jfull}
\end{equation}

Now substitute~\eqref{eq:Jfull} into~\eqref{eq:poissonA}:

\begin{equation}
  \nabla^2 \vA = -\muZ \sig\left(-\nabla\phi - \Der{\vA}{t} + \vv \times (\nabla \times \vA)\right)
\end{equation}

\begin{equation}
  \nabla^2 \vA = \muZ \sig \nabla\phi
               + \muZ \sig \Der{\vA}{t}
               - \muZ \sig \,\vv \times (\nabla \times \vA)
  \label{eq:Afull}
\end{equation}

\subsection{Eliminating the Scalar Potential}

Taking the divergence of both sides of~\eqref{eq:Jfull} and applying $\nabla \cdot \vJ = 0$
(charge conservation in the conductor) and the Coulomb gauge $\nabla \cdot \vA = 0$:

\begin{equation}
  0 = \sig\left(-\nabla^2\phi - \nabla \cdot \Der{\vA}{t}
      + \nabla \cdot (\vv \times (\nabla \times \vA))\right)
\end{equation}

For uniform velocity $\vv = v\,\xhat$ (no spatial variation in $\vv$):

\begin{equation}
  \nabla \cdot (\vv \times \vB) = \vB \cdot (\nabla \times \vv) - \vv \cdot (\nabla \times \vB)
  = 0 - v\,(\nabla \times \vB)_x = -v\,(\muZ J_x)
\end{equation}

In the geometry of interest the rotor current flows purely in $\yhat$ (shown in
Section~\ref{sec:currentdir}), so $J_x = 0$ and:

\begin{equation}
  \nabla^2\phi = -\frac{\partial}{\partial t}\underbrace{(\nabla \cdot \vA)}_{=0} = 0
\end{equation}

With appropriate boundary conditions (no charge accumulation on the rotor surfaces),
$\phi = 0$ inside the conductor. Equation~\eqref{eq:Afull} then reduces to:

\begin{tcolorbox}[colback=yellow!10, colframe=orange!60,
  title=\textbf{Convective Magnetic Diffusion Equation (CMDE)}]
\begin{equation}
  \boxed{
    \nabla^2 \vA = \muZ \sig \Der{\vA}{t} + \muZ \sig\,(\vv \cdot \nabla)\vA
  }
  \label{eq:cmde}
\end{equation}
\end{tcolorbox}

\noindent where we used the identity
$\vv \times (\nabla \times \vA) = \nabla(\vv \cdot \vA) - (\vv \cdot \nabla)\vA$
and noted $\nabla(\vv \cdot \vA) = \nabla(v A_x)$ contributes only to the $\phi$
equation (already set to zero).

\subsection{Physical Interpretation of the Two RHS Terms}

\begin{enumerate}[label=\textbf{\arabic*.}]
  \item \textbf{Temporal diffusion:} $\muZ\sig\,\partial\vA/\partial t$ drives the
        magnetic skin effect. At high slip frequency $\oms$, the penetration depth
        $\delta = \sqrt{2/(\muZ\sig\oms)}$ shrinks, confining eddy currents to a thin
        surface layer.
  \item \textbf{Convective distortion:} $\muZ\sig\,(\vv\cdot\nabla)\vA$ represents
        the dragging of field lines by rotor motion. At high rotor speed, the field
        is compressed on the leading edge and rarefied on the trailing edge — the
        physical origin of the LIM end-effect efficiency penalty. At low velocity this
        term is negligible relative to the diffusion term.
\end{enumerate}

\subsection{Direction of Induced Current}
\label{sec:currentdir}

The stator field has $\vB_s = B_0\cos(kx - \omega t)\,\zhat$. For a rotor moving
in $\xhat$, the motional EMF term $\vv \times \vB$ points in $\xhat \times \zhat = -\yhat$.
Faraday induction from a $z$-directed time-varying field also drives currents in the
$x$-$y$ plane. Since we consider a thin rotor (dimension in $z$ much smaller than in
$x$ and $y$), current loops close in $\yhat$ across the rotor width. Therefore
$\vA = A_z(x,y,t)\,\zhat$ and $\vJ = J_y(x,y,t)\,\yhat$.


% ============================================================
\section{Nondimensionalisation and Governing Dimensionless Groups}
\label{sec:nondim}

\subsection{Scaling Variables}

Let the characteristic length scale be the pole pitch $\taup$, the characteristic time scale
be $1/\oms$ (one slip cycle), and the characteristic vector potential scale be $B_0\taup$.
Define the dimensionless variables:

\begin{equation}
  \tilde{x} = \frac{x}{\taup}, \quad
  \tilde{y} = \frac{y}{\taup}, \quad
  \tilde{t} = \oms t, \quad
  \tilde{A}_z = \frac{A_z}{B_0\taup}
\end{equation}

\subsection{Dimensionless CMDE}

Substituting into equation~\eqref{eq:cmde} (retaining the $z$-component only, since
$\vA = A_z\,\zhat$):

\begin{align}
  \frac{1}{\taup^2}\left(\DDer{\tilde{A}_z}{\tilde{x}} + \DDer{\tilde{A}_z}{\tilde{y}}\right)
  &= \muZ\sig\,\oms\,\frac{1}{\taup^2}\Der{\tilde{A}_z}{\tilde{t}}
   + \muZ\sig\,v\,\frac{1}{\taup^2}\Der{\tilde{A}_z}{\tilde{x}}
\end{align}

Dividing through by $1/\taup^2$:

\begin{equation}
  \DDer{\tilde{A}_z}{\tilde{x}} + \DDer{\tilde{A}_z}{\tilde{y}}
  = \underbrace{\muZ\sig\oms\taup^2}_{\displaystyle S}
    \Der{\tilde{A}_z}{\tilde{t}}
  + \underbrace{\muZ\sig v\taup}_{\displaystyle \Rm}
    \Der{\tilde{A}_z}{\tilde{x}}
  \label{eq:cmde_nd}
\end{equation}

\begin{tcolorbox}[colback=green!5, colframe=green!50,
  title=\textbf{The Two Governing Dimensionless Groups}]
\begin{align}
  S &= \muZ\sig\oms\taup^2
     && \text{(Quality factor / dimensionless skin depth parameter)} \label{eq:Sdef}\\
  \Rm &= \muZ\sig v\taup
     && \text{(Magnetic Reynolds number)} \label{eq:Rmdef}
\end{align}
\end{tcolorbox}

\begin{remark}
  The ratio $S/\Rm = \oms\taup/v = (\omega - kv)\cdot(\taup/v) = \pi s$ where $s$
  is the conventional slip. This recovers the familiar induction-machine result and
  confirms that the nondimensionalisation is physically consistent.
\end{remark}

\begin{remark}
  \textbf{Equivalence principle:} Equation~\eqref{eq:cmde_nd} shows that two PCB-LIM
  systems with equal $S$ and $\Rm$ are physically equivalent — they produce identical
  dimensionless thrust $\Fstar$. This is the mathematical foundation for the proposed
  scaling law: one curve of $\Fstar$ vs.\ $S$ (at fixed $\Rm \ll 1$) characterises
  all low-speed planar LIM geometries.
\end{remark}

% ============================================================
\section{Magnetic Skin Depth and Regime Classification}
\label{sec:skindepth}

\subsection{Skin Depth Derivation}

Consider a planar conductor occupying $y \leq 0$ (rotor surface at $y=0$), driven by a
surface current oscillating at slip frequency $\oms$. In the stationary rotor frame
($v=0$, $\Rm = 0$), the CMDE~\eqref{eq:cmde} reduces to:

\begin{equation}
  \DDer{A_z}{y} = \muZ\sig\oms \jmath A_z
\end{equation}

where we use the complex phasor $A_z(y,t) = \hat{A}(y)\,e^{\jmath\oms t}$. The general
solution decaying into the conductor ($y \to -\infty$) is:

\begin{equation}
  \hat{A}(y) = \hat{A}_0 \, e^{\kappa y}, \quad
  \kappa = \sqrt{\jmath\muZ\sig\oms} = \frac{1+\jmath}{\delta}
\end{equation}

where the \emph{magnetic skin depth} is:

\begin{tcolorbox}[colback=yellow!10, colframe=orange!60]
\begin{equation}
  \boxed{\delta = \sqrt{\frac{2}{\muZ\sig\oms}}}
  \label{eq:skindepth}
\end{equation}
\end{tcolorbox}

The field amplitude decays as $e^{-y/\delta}$ and is rotated in phase by $e^{-\jmath y/\delta}$
(one radian per skin depth of penetration).

\subsection{Thin-Sheet vs.\ Skin-Effect Regimes}

Let $d$ denote the rotor thickness. Two physically distinct regimes arise:

\begin{description}
  \item[Thin-sheet regime ($d \ll \delta$):] The field penetrates the entire rotor
        uniformly. Current density is approximately constant across $y$. Thrust scales
        as $F \propto \sig d^2 \oms$ (confirmed by Theorem~\ref{thm:volcancellation}).

  \item[Skin-effect regime ($d \gg \delta$):] Eddy currents are confined to a surface
        layer of thickness $\delta$. Thrust saturates and becomes independent of $d$.
        Thrust scales as $F \propto \sig^{1/2} \oms^{-1/2}$ from the Maxwell Stress
        Tensor result (eq.~\eqref{eq:Fstar_final}).
\end{description}

The transition occurs when $\delta = d$:

\begin{equation}
  \frac{2}{\muZ\sig\oms^{\mathrm{tr}}} = d^2
  \quad \Longrightarrow \quad
  \oms^{\mathrm{tr}} = \frac{2}{\muZ\sig d^2}
  \label{eq:omtr}
\end{equation}

\begin{remark}
  For 7075-T6 aluminium ($\sig = 1.96\times10^7\,\si{S/m}$) and rotor thickness
  $d = 3\,\si{mm}$:
  \[
    \oms^{\mathrm{tr}}
    = \frac{2}{(4\pi\times10^{-7})(1.96\times10^7)(3\times10^{-3})^2}
    \approx 900\,\si{rad/s}
    \quad \Longrightarrow \quad
    f^{\mathrm{tr}} \approx 143\,\si{Hz}
  \]
  The proposed parametric experiment ($\oms/2\pi \in \{50, 150, 400\}\,\si{Hz}$)
  deliberately straddles this transition.
\end{remark}


% ============================================================
\section{Boundary Value Problem in the Rotor Half-Space}
\label{sec:bvp}

\subsection{Setup}

We solve the CMDE in the rotor rest frame. Place the stator surface at $y=0$
and the rotor occupying $-d \leq y \leq 0$. The stator imposes a boundary
condition at $y = 0^+$ (just above the rotor surface):

\begin{equation}
  A_z\big|_{y=0} = \hat{A}_0\, e^{\jmath(\pi x/\taup - \omega t)}
  \label{eq:bc_stator}
\end{equation}

corresponding to a sinusoidal travelling-wave surface current
$K_z = K_0\cos(\pi x/\taup - \omega t)$. In the rotor rest frame, this boundary
condition oscillates at the slip frequency $\oms = \omega - kv$, so:

\begin{equation}
  A_z\big|_{y=0} = \hat{A}_0\, e^{\jmath(\pi x/\taup - \oms t)}
\end{equation}

Seek a solution of the form:
\begin{equation}
  A_z(x,y,t) = \mathrm{Re}\left[\hat{A}(y)\,e^{\jmath(\pi x/\taup - \oms t)}\right]
  \label{eq:ansatz}
\end{equation}

\subsection{ODE for the $y$-Dependence}

Substituting ansatz~\eqref{eq:ansatz} into the CMDE~\eqref{eq:cmde} with $v=0$
in the rotor frame:

\begin{align}
  \left[-\left(\frac{\pi}{\taup}\right)^2 \hat{A} + \DDer{\hat{A}}{y}\right]
  e^{\jmath(\pi x/\taup - \oms t)}
  &= \muZ\sig(-\jmath\oms)\hat{A}\,
     e^{\jmath(\pi x/\taup - \oms t)}
\end{align}

Cancelling the common exponential factor:

\begin{equation}
  \DDer{\hat{A}}{y} = \left[\left(\frac{\pi}{\taup}\right)^2 - \jmath\muZ\sig\oms\right]\hat{A}
  \equiv \kappa^2\hat{A}
  \label{eq:ode_y}
\end{equation}

where:
\begin{equation}
  \kappa^2 = \left(\frac{\pi}{\taup}\right)^2\left[1 - \jmath S\right],
  \quad S = \frac{\muZ\sig\oms\taup^2}{\pi^2}
  \label{eq:kappa}
\end{equation}

Note that $S$ here is the same quality factor defined in~\eqref{eq:Sdef} divided by
$\pi^2$; this is a common alternative normalisation. We adopt the convention that
$S \equiv \muZ\sig\oms\taup^2/\pi^2$ throughout the force calculation.

\subsection{Solution of the ODE}

The general solution to~\eqref{eq:ode_y} is:
\begin{equation}
  \hat{A}(y) = C_1 e^{\kappa y} + C_2 e^{-\kappa y}
\end{equation}

Boundary conditions:
\begin{itemize}
  \item At $y = 0$: $\hat{A}(0) = \hat{A}_0$ (stator BC~\eqref{eq:bc_stator}).
  \item At $y = -d$: $\hat{A}(-d) = \hat{A}_0\,g(d)$ where $g$ depends on the
        regime. For the thin-sheet approximation ($d \ll \delta$, i.e.\ $|\kappa|d \ll 1$):
        $g \approx 1$ and we can use $\hat{A}(-d) \approx \hat{A}_0$, giving
        equal amplitudes at both surfaces (uniform penetration).
\end{itemize}

\noindent\textbf{Thin-sheet simplification ($|\kappa|d \ll 1$):} Taylor-expanding:
\begin{equation}
  e^{\pm\kappa y} \approx 1 \pm \kappa y + \tfrac{1}{2}(\kappa y)^2 + \ldots
\end{equation}
To leading order, $\hat{A}(y) \approx \hat{A}_0 = \text{const}$. The first
correction to the current density is:

\begin{align}
  \hat{J}_y &= -\frac{\sig}{\jmath}\cdot\frac{\partial(\jmath\oms\hat{A})}{\partial \cdot}
              \approx \sig\oms\hat{A}_0 \cdot (s \text{ dependence})
\end{align}

The full derivation via the thin-sheet integral gives the eddy current density
(uniform in $y$):

\begin{tcolorbox}[colback=yellow!10, colframe=orange!60]
\begin{equation}
  \boxed{J_y = \sig\, s\, v_s\, B_0 \cos(kx - \omega t)}
  \label{eq:Jy}
\end{equation}
\end{tcolorbox}

\noindent\textit{Derivation of~\eqref{eq:Jy}:} In the rotor rest frame, the
motional EMF per unit length is $\mathcal{E} = v_{\mathrm{rel}}B_0 = sv_sB_0$
(the slip velocity times the field). For a thin-sheet of thickness $d$ and
conductivity $\sig$, the current density is
$J_y = \sig\,\mathcal{E} \cdot f_{\mathrm{wave}}(x,t)$, confirming~\eqref{eq:Jy}.

% ============================================================
\section{Lorentz Force and Time-Averaged Thrust}
\label{sec:lorentz}

\subsection{Instantaneous Volumetric Force Density}

The Lorentz force density on the rotor is:
\begin{equation}
  \mathbf{f} = \vJ \times \vB = J_y\,\yhat \times B_z\,\zhat = J_y B_z\,\xhat
\end{equation}

Substituting~\eqref{eq:Jy} and $B_z = B_0\cos(kx-\omega t)$:

\begin{align}
  f_x &= \left[\sig\, s\, v_s\, B_0\cos(kx-\omega t)\right]
         \cdot \left[B_0\cos(kx-\omega t)\right] \nonumber\\
      &= \sig\, s\, v_s\, B_0^2\cos^2(kx-\omega t)
  \label{eq:fx_inst}
\end{align}

\subsection{Time-Averaging}

The physical quantity governing continuous thrust is the time average of~\eqref{eq:fx_inst}
over one period $T = 2\pi/\omega$. Since $\cos^2\theta$ averages to $1/2$ over any
full cycle:

\begin{equation}
  \avg{f_x} = \frac{1}{T}\int_0^T \sig\, s\, v_s\, B_0^2\cos^2(kx-\omega t)\,dt
  = \frac{1}{2}\sig\, s\, v_s\, B_0^2
  \label{eq:fx_avg}
\end{equation}

\subsection{Total Thrust on the Rotor}

The total time-averaged propulsive force on a rotor of volume $V = \ell\, w\, d$
(length $\times$ width $\times$ thickness) is:

\begin{equation}
  F_x = \avg{f_x} \cdot V = \frac{1}{2}\sig\, V\, s\, v_s\, B_0^2
  \label{eq:Fx_total}
\end{equation}

% ============================================================
\section{Maxwell Stress Tensor Formulation}
\label{sec:mst}

\subsection{Definition}

The Maxwell Stress Tensor provides an independent, rigorous derivation of the
electromagnetic force on a body by integrating a field stress over its bounding
surface. For a non-magnetic medium ($\mu = \muZ$):

\begin{equation}
  T_{ij} = \frac{1}{\muZ}\left(B_i B_j - \frac{1}{2}\delta_{ij}B^2\right)
  \label{eq:MST}
\end{equation}

The time-averaged force per unit area (traction) on a surface with outward
normal $\nhat$ is:

\begin{equation}
  \avg{f_i} = \avg{T_{ij}}\,n_j
\end{equation}

\subsection{Thrust from the Air-Gap Tangential Stress}

The propulsive ($x$-directed) thrust is obtained by integrating the shear stress
$\avg{T_{xy}}$ over the air-gap surface at $y = 0$ (stator surface), with $\nhat = \yhat$:

\begin{align}
  \avg{F_x} &= w \int_0^{2\taup} \avg{T_{xy}}\big|_{y=0}\,dx
             = \frac{w}{\muZ}\int_0^{2\taup}\avg{B_x B_y}\big|_{y=0}\,dx
  \label{eq:Fx_MST}
\end{align}

where $w$ is the rotor width and the integral spans two pole pitches (one full
spatial period).

\subsection{Analytical Evaluation at Low $\Rm$}

In the limit $\Rm \ll 1$ (low rotor velocity relative to the Alfv\'en speed), the
field inside the rotor is the superposition of the applied stator field and the
field due to induced currents. Using the complex exponential solution from
Section~\ref{sec:bvp}, one can show (see Appendix~\ref{app:mst_calc}) that:

\begin{equation}
  \avg{B_x B_y}\big|_{y=0} = \frac{B_0^2}{2}\cdot\frac{S}{1+S^2}
\end{equation}

Substituting into~\eqref{eq:Fx_MST} and integrating over $2\taup$:

\begin{tcolorbox}[colback=yellow!10, colframe=orange!60,
  title=\textbf{Closed-Form Thrust Scaling Law}]
\begin{equation}
  \boxed{
    \avg{F_x} = \frac{B_0^2\,\taup\,w}{2\muZ} \cdot \frac{S}{1+S^2}
  }
  \label{eq:Fx_MST_final}
\end{equation}
\end{tcolorbox}

Defining the dimensionless thrust:

\begin{equation}
  \Fstar \equiv \frac{\avg{F_x}\,\muZ}{B_0^2\,\taup\,w}
  \quad \Longrightarrow \quad
  \boxed{\Fstar = \frac{S}{1+S^2}}
  \label{eq:Fstar_final}
\end{equation}

This is the central falsifiable prediction of the thesis.

% ============================================================
\section{Optimal Slip Frequency}
\label{sec:optslip}

\begin{proposition}[Optimal Slip Frequency]
\label{prop:optslip}
  Under a sinusoidal travelling-wave stator field and the boundary-value problem
  of Section~\ref{sec:bvp}, the dimensionless thrust $\Fstar(S)$ attains its maximum
  at:
  \begin{equation}
    S^* = 1 \quad \Longleftrightarrow \quad
    \omstar = \frac{\pi^2}{\muZ\sig\taup^2}
    \label{eq:omstar}
  \end{equation}
  At the optimum, $\Fstar(S^*) = 1/2$ and the optimal skin depth satisfies
  $\delta^* = \sqrt{2}\,\taup/\pi$.
\end{proposition}

\begin{proof}
  Differentiate $\Fstar = S/(1+S^2)$ with respect to $S$:
  \begin{align}
    \frac{d\Fstar}{dS}
    &= \frac{(1+S^2) - S\cdot 2S}{(1+S^2)^2}
     = \frac{1 - S^2}{(1+S^2)^2}
    \label{eq:dFstar}
  \end{align}
  Setting to zero: $1 - S^2 = 0 \Rightarrow S^* = 1$ (taking the positive root).
  This is a maximum since $d^2\Fstar/dS^2\big|_{S=1} = -1/2 < 0$.

  Since $S = \muZ\sig\oms\taup^2/\pi^2$, setting $S^* = 1$ gives:
  \begin{equation}
    \muZ\sig\omstar\taup^2/\pi^2 = 1
    \quad \Longrightarrow \quad
    \omstar = \frac{\pi^2}{\muZ\sig\taup^2} = \frac{\pi^2\eta}{\taup^2}
  \end{equation}
  where $\eta = 1/(\muZ\sig)$ is the magnetic diffusivity.

  At $S^* = 1$: $\delta^* = \sqrt{2/(\muZ\sig\omstar)} = \sqrt{2\taup^2/\pi^2}
  = \sqrt{2}\,\taup/\pi$.
\end{proof}

\begin{remark}[Scaling Consequences]
  Proposition~\ref{prop:optslip} makes two experimentally testable predictions:
  \begin{enumerate}
    \item Doubling $\taup$ reduces $\omstar$ by a factor of four ($\omstar \propto \taup^{-2}$).
    \item Replacing 7075-T6 aluminium with 6061-T6 (a $\approx 15\%$ reduction in $\sig$)
          shifts $\omstar$ upward by $15\%$.
  \end{enumerate}
  Both will be directly tested in the 18-point factorial experiment (Phase 5).
\end{remark}


% ============================================================
\section{Volume Cancellation Theorem}
\label{sec:volcancellation}

\begin{theorem}[Volume Cancellation]
\label{thm:volcancellation}
  In the thin-sheet limit ($d \ll \delta$), the rotor equation of motion reduces to
  a first-order ODE in velocity $v(t)$ that is \emph{independent of rotor plan-area}
  $A = \ell\,w$. The acceleration depends on rotor thickness $d$ only through the
  ratio $\sig\,d / \rho$.
\end{theorem}

\begin{proof}
  \textbf{Step 1: Induced EMF.}
  In the thin-sheet limit, the magnetic flux linking the rotor cross-section changes
  at a rate proportional to the relative velocity between rotor and wave:
  \begin{equation}
    \mathcal{E} = (v_s - v)\,B_0 = s\,v_s\,B_0
  \end{equation}
  per unit length along $\yhat$ (the rotor width direction).

  \textbf{Step 2: Short-circuit current density.}
  The thin sheet of thickness $d$ and conductivity $\sig$ has sheet resistance
  $R_\square = 1/(\sig d)$. The short-circuit current density (current per unit
  width in $\yhat$) is:
  \begin{equation}
    K_{\mathrm{ind}} = \sig\,d\,\mathcal{E} = \sig\,d\,s\,v_s\,B_0
  \end{equation}
  and the volumetric current density is $J_y = K_{\mathrm{ind}}/d = \sig\,s\,v_s\,B_0$,
  confirming~\eqref{eq:Jy} for the thin-sheet case.

  \textbf{Step 3: Lorentz force.}
  The force density is $f_x = J_y B_0 = \sig\,s\,v_s\,B_0^2$. The time-averaged
  force density is $\avg{f_x} = \tfrac{1}{2}\sig\,s\,v_s\,B_0^2$ (Section~\ref{sec:lorentz}).
  Total force:
  \begin{equation}
    F_x = \avg{f_x}\cdot V = \frac{1}{2}\sig\,\ell\,w\,d\cdot s\,v_s\,B_0^2
    \label{eq:Fx_vol}
  \end{equation}

  \textbf{Step 4: Newton's Second Law.}
  The rotor mass is $m = \rho\,\ell\,w\,d = \rho\,V$. Newton's law gives:
  \begin{equation}
    m\,\dot{v} = F_x
    \quad \Longrightarrow \quad
    \rho\,\cancel{\ell\,w\,d}\,\dot{v}
    = \frac{1}{2}\sig\,\cancel{\ell\,w\,d}\cdot s\,v_s\,B_0^2
  \end{equation}

  The volume $V = \ell\,w\,d$ cancels \emph{exactly} from both sides:

  \begin{tcolorbox}[colback=yellow!10, colframe=orange!60]
  \begin{equation}
    \boxed{
      \rho\,\dot{v} = \frac{1}{2}\sig\,(v_s - v)\,B_0^2
    }
    \label{eq:newton_reduced}
  \end{equation}
  \end{tcolorbox}

  \textbf{Step 5: Independence of plan-area.}
  Equation~\eqref{eq:newton_reduced} contains only $\rho$, $\sig$, $B_0$, $v_s$,
  and $v$. The geometric quantities $\ell$, $w$, $d$ do not appear. The acceleration
  is therefore independent of rotor length, width, and thickness in the thin-sheet
  limit.
  
  Note: the rotor thickness $d$ \emph{does} appear implicitly via $\sig$: if $d$
  exceeds the skin depth, the effective $\sig_{\mathrm{eff}} = \sig\delta/d$ (skin-effect
  saturation) must be used, restoring $d$ dependence in the thick-sheet regime.
  In the thin-sheet limit ($d \ll \delta$) stated by the theorem, $\sig_{\mathrm{eff}}
  = \sig$ and the result holds exactly. \qed
\end{proof}

\begin{remark}[Experimental Implication]
  Theorem~\ref{thm:volcancellation} implies that rotor mass need not be controlled
  between experimental runs — only the material ratio $\sig/\rho$ matters.
  For 7075-T6 and 6061-T6 aluminium alloys:
  \begin{align}
    \left(\frac{\sig}{\rho}\right)_{7075}
      &= \frac{1.96\times10^7}{2810} = 6974\,\si{S\cdot m^2 / kg} \\
    \left(\frac{\sig}{\rho}\right)_{6061}
      &= \frac{2.33\times10^7}{2700} = 8630\,\si{S\cdot m^2 / kg}
  \end{align}
  This $\approx 24\%$ difference in $\sig/\rho$ is comfortably resolved by the
  measurement system (combined uncertainty $< 4\%$ on inferred thrust,
  Table~6 of the proposal).
\end{remark}


% ============================================================
\section{Kinematic ODE and Exact Solution}
\label{sec:kinematics}

\subsection{Derivation of the First-Order ODE}

Rearranging~\eqref{eq:newton_reduced} and substituting $s = (v_s - v)/v_s$:

\begin{equation}
  \frac{dv}{dt} = \frac{\sig B_0^2}{2\rho}(v_s - v)
  \label{eq:dvdt}
\end{equation}

Define the electrodynamic time constant:
\begin{equation}
  \alpha \equiv \frac{\sig B_0^2}{2\rho}
  \label{eq:alpha}
\end{equation}

Then~\eqref{eq:dvdt} is a linear first-order ODE:
\begin{equation}
  \frac{dv}{dt} + \alpha\,v = \alpha\,v_s
  \label{eq:ode_v}
\end{equation}

\subsection{Exact Analytical Solution}

The integrating factor is $e^{\alpha t}$. Multiplying both sides:

\begin{equation}
  \frac{d}{dt}\left(e^{\alpha t} v\right) = \alpha\,v_s\,e^{\alpha t}
\end{equation}

Integrating from $t=0$ to $t$ with initial condition $v(0) = 0$:

\begin{equation}
  e^{\alpha t} v(t) - 0 = \alpha\,v_s \int_0^t e^{\alpha\tau}\,d\tau
  = v_s\left(e^{\alpha t} - 1\right)
\end{equation}

\begin{tcolorbox}[colback=green!5, colframe=green!50,
  title=\textbf{Exact Velocity Solution}]
\begin{equation}
  \boxed{v(t) = v_s\left(1 - e^{-\alpha t}\right)}
  \label{eq:vt}
\end{equation}
\end{tcolorbox}

Integrating once more for position ($x(0) = 0$):

\begin{equation}
  x(t) = \int_0^t v(\tau)\,d\tau
        = v_s\left[t + \frac{1}{\alpha}e^{-\alpha t} - \frac{1}{\alpha}\right]
        = v_s\left(t - \frac{1-e^{-\alpha t}}{\alpha}\right)
  \label{eq:xt}
\end{equation}

\subsection{Characteristic Behaviour}

\begin{itemize}
  \item \textbf{Early time ($t \ll 1/\alpha$):} $v(t) \approx \alpha\,v_s\,t$, i.e.\
        constant acceleration $a_0 = \alpha\,v_s = \sig\,B_0^2\,v_s/(2\rho)$.
  \item \textbf{Late time ($t \gg 1/\alpha$):} $v(t) \to v_s$ exponentially; the
        projectile asymptotically approaches the synchronous wave speed.
  \item \textbf{Time constant:} $\tau = 1/\alpha = 2\rho/(\sig B_0^2)$.
        For 7075-T6 with $B_0 = 10\,\si{mT}$:
        \[
          \tau = \frac{2 \times 2810}{1.96\times10^7 \times (10^{-2})^2}
               \approx 2.87\,\si{s}
        \]
        This matches the ODE simulation output in \texttt{kinematic\_solver.py}
        (Figure~1 of the proposal).
\end{itemize}

\subsection{Required Stator Length}

The PCB stator must be at least as long as the distance $x(t^*)$ at which the
projectile reaches the exit condition. For a target exit velocity $v_{\mathrm{exit}}$:

\begin{equation}
  t^* = -\frac{1}{\alpha}\ln\!\left(1 - \frac{v_{\mathrm{exit}}}{v_s}\right)
\end{equation}

\begin{equation}
  L_{\mathrm{stator}} = x(t^*)
  = v_s\left(t^* - \frac{1 - e^{-\alpha t^*}}{\alpha}\right)
  \label{eq:Lstator}
\end{equation}

Equation~\eqref{eq:Lstator} is evaluated numerically in
\texttt{simulations/kinematic\_solver.py} to dictate the required PCB stator length
during the hardware design phase (Phase 3).


% ============================================================
\section{Summary of Key Equations}
\label{sec:summary}

\begin{table}[h!]
  \centering
  \renewcommand{\arraystretch}{1.5}
  \begin{tabular}{lll}
    \toprule
    \textbf{Quantity} & \textbf{Expression} & \textbf{Eq.}\\
    \midrule
    CMDE (governing PDE) &
      $\nabla^2\vA = \muZ\sig\partial\vA/\partial t + \muZ\sig(\vv\cdot\nabla)\vA$ &
      \eqref{eq:cmde}\\
    Quality factor &
      $S = \muZ\sig\oms\taup^2/\pi^2$ &
      \eqref{eq:kappa}\\
    Magnetic Reynolds number &
      $\Rm = \muZ\sig v\taup$ &
      \eqref{eq:Rmdef}\\
    Skin depth &
      $\delta = \sqrt{2/(\muZ\sig\oms)}$ &
      \eqref{eq:skindepth}\\
    Transition frequency &
      $\oms^{\mathrm{tr}} = 2/(\muZ\sig d^2)$ &
      \eqref{eq:omtr}\\
    Dimensionless thrust &
      $\Fstar = S/(1+S^2)$ &
      \eqref{eq:Fstar_final}\\
    Optimal slip frequency &
      $\omstar = \pi^2/(\muZ\sig\taup^2)$ &
      \eqref{eq:omstar}\\
    Volume-cancelled EOM &
      $\rho\dot{v} = \tfrac{1}{2}\sig(v_s-v)B_0^2$ &
      \eqref{eq:newton_reduced}\\
    Exact velocity solution &
      $v(t) = v_s(1 - e^{-\alpha t})$, $\alpha = \sig B_0^2/(2\rho)$ &
      \eqref{eq:vt}\\
    \bottomrule
  \end{tabular}
  \caption{Complete equation reference for the PCB-LIM analytical model.}
  \label{tab:equations}
\end{table}


% ============================================================
\appendix
\section{Maxwell Stress Tensor Calculation}
\label{app:mst_calc}

\subsection{Field Components in the Air Gap}

The total field at $y = 0^+$ (stator surface, just above the rotor) is the superposition
of the applied stator field and the reaction field due to induced currents:

\begin{align}
  B_x &= -\Der{A_z}{y}\bigg|_{y=0} \\
  B_y &= \frac{\jmath\pi}{\taup}\hat{A}(0)
\end{align}

From the BVP solution~\eqref{eq:ode_y} with $\kappa^2 = (\pi/\taup)^2(1-\jmath S)$:

\begin{align}
  \hat{A}(0) &= \hat{A}_0 \\
  \Der{\hat{A}}{y}\bigg|_{y=0} &= \kappa\hat{A}_0\,\tanh(\kappa d)
                                 \approx \kappa^2 d\,\hat{A}_0 \quad (|\kappa|d \ll 1)
\end{align}

\subsection{Time-Averaged Stress}

Using $\avg{\mathrm{Re}(\hat{B}_x e^{-\jmath\oms t})\,\mathrm{Re}(\hat{B}_y e^{-\jmath\oms t})}
= \tfrac{1}{2}\mathrm{Re}(\hat{B}_x\hat{B}_y^*)$:

\begin{align}
  \avg{B_x B_y}\big|_{y=0}
    &= \frac{1}{2}\mathrm{Re}\!\left[-\kappa^2 d\,|\hat{A}_0|^2 \cdot
       \frac{-\jmath\pi}{\taup}\right] \nonumber\\
    &= \frac{1}{2}\mathrm{Re}\!\left[\jmath\frac{\pi}{\taup}
       \left(\frac{\pi}{\taup}\right)^2(1-\jmath S)\,d\,|\hat{A}_0|^2\right] \nonumber\\
    &= \frac{|\hat{A}_0|^2\,d}{2}\left(\frac{\pi}{\taup}\right)^3\mathrm{Re}[\jmath(1-\jmath S)] \nonumber\\
    &= \frac{|\hat{A}_0|^2\,d}{2}\left(\frac{\pi}{\taup}\right)^3 S
\end{align}

Since $|\hat{A}_0| = B_0\taup/\pi$ (from the relation $B_y = -\jmath k A_z$ at the surface
with $k = \pi/\taup$), and collecting factors:

\begin{equation}
  \avg{B_x B_y}\big|_{y=0}
  = \frac{B_0^2}{2}\cdot\frac{S}{1+S^2}
\end{equation}

Integrating over $2\taup$ and multiplying by $w/\muZ$ gives~\eqref{eq:Fx_MST_final}.
$\square$


\section{Parameter Values for 7075-T6 and 6061-T6 Aluminium}
\label{app:params}

\begin{table}[h!]
  \centering
  \renewcommand{\arraystretch}{1.3}
  \begin{tabular}{lccc}
    \toprule
    \textbf{Property} & \textbf{Symbol} & \textbf{7075-T6} & \textbf{6061-T6}\\
    \midrule
    Electrical conductivity & $\sig$ & $1.96\times10^7\,\si{S/m}$ & $2.33\times10^7\,\si{S/m}$\\
    Mass density & $\rho$ & $2810\,\si{kg/m^3}$ & $2700\,\si{kg/m^3}$\\
    Relative permeability & $\mu_r$ & 1.0 (non-magnetic) & 1.0 (non-magnetic)\\
    $\sig/\rho$ ratio & & $6974\,\si{S\cdot m^2/kg}$ & $8630\,\si{S\cdot m^2/kg}$\\
    $\omstar$ ($\taup = 10\,\si{mm}$) & & $\approx 508\,\si{rad/s}$ & $\approx 427\,\si{rad/s}$\\
    $f^* = \omstar/2\pi$ & & $\approx 81\,\si{Hz}$ & $\approx 68\,\si{Hz}$\\
    \bottomrule
  \end{tabular}
  \caption{Material parameters for the two rotor alloys. $\omstar$ computed from~\eqref{eq:omstar}.}
\end{table}


\section{Numerical Parameters for the ODE Solver}
\label{app:ode_params}

\begin{table}[h!]
  \centering
  \renewcommand{\arraystretch}{1.3}
  \begin{tabular}{lccc}
    \toprule
    \textbf{Parameter} & \textbf{Symbol} & \textbf{Value} & \textbf{Unit}\\
    \midrule
    Stator frequency & $f$ & 100 & Hz\\
    Pole pitch & $\taup$ & 15 & mm\\
    Peak flux density & $B_0$ & 10 & mT\\
    Synchronous speed & $v_s = 2f\taup$ & 3.0 & m/s\\
    Electrodynamic time constant & $\tau = 1/\alpha$ & 2.87 & s\\
    Simulation duration & $t_{\max}$ & 1.5 & s\\
    Time steps & & 1000 & —\\
    \bottomrule
  \end{tabular}
  \caption{Parameters used in \texttt{simulations/kinematic\_solver.py}.}
\end{table}

\end{document}